% Source: https://artofproblemsolving.com/community/c3h1433925_ellipses
% Stable identity: aops-topic-1433925-q01
The point reflection in $M$ (rotation by $180^\circ$ about $M$) sends $P$ to $P'$, keeps distances, and swaps $F_1$ and $F_2$ (because $M$ is the midpoint of $F_1F_2$). So $P'F_2=PF_1$ and $P'F_1=PF_2$, and
\[
P'F_1+P'F_2=PF_2+PF_1=2a .
\]
Hence $P'$ is on the ellipse. (In the figure the red dashed path $F_1P'F_2$ is the image of the blue path $F_1PF_2$.)

\begin{asy}
usepackage("amsmath");
size(6cm);
defaultpen(fontsize(10pt));
real a=3,c=2,bb=sqrt(a^2-c^2);
pair F1=(-c,0),F2=(c,0),M=(0,0);
draw(ellipse(M,a,bb),linewidth(0.9));
pair P=(a*cos(1.0),bb*sin(1.0)), Q=-P;
draw(F1--P--F2,blue); draw(F1--Q--F2,red+dashed);
draw(P--Q,gray+dashed);
dot(F1);dot(F2);dot(M);dot(P);dot(Q);
label("$F_1$",F1,W); label("$F_2$",F2,E); label("$M$",M,SE); label("$P$",P,NE); label("$P'$",Q,SW);
label("\mbox{}",(min(currentpicture,user=true)+max(currentpicture,user=true))/2,black); // keeps the text after the figure black when LaTeX typesets the labels inline
\end{asy}

The \emph{centre} of the ellipse is its centre of symmetry, the point $M$: every chord through $M$ is bisected at $M$. So it is not an extra definition that the centre is equidistant from the foci; it is a consequence: the reflection argument shows that $M$, the midpoint of the foci, is a centre of symmetry. Moreover, it is the only one: if $Z_1$ and $Z_2$ were two centres of symmetry, doing the reflection in $Z_1$ and then the reflection in $Z_2$ would be the translation by the vector $2\,\overrightarrow{Z_1Z_2}$, and it would map the ellipse onto itself; repeating it would move points of the ellipse arbitrarily far away, which is impossible for a bounded figure unless $Z_1=Z_2$. Thus the centre is the midpoint of $F_1F_2$, at distance $\dfrac{F_1F_2}{2}$ from each focus. $\blacksquare$
